\section{Conclusion}

In this paper, we present \method, a discrete flow-matching VLA framework that enables iterative full-sequence action refinement through token-level velocity modeling. Our \method~can revisit and correct previously generated action tokens, improving action consistency for robotic manipulation.
We further study two complementary velocity-field constructions and integrate a practical two-stage decoder. Extensive experiments demonstrate that \method~consistently delivers strong action quality while maintaining high inference efficiency.


% In this paper, we presented \method, a discrete flow matching framework for the VLA model that enables iterative full-sequence action refinement. Unlike autoregressive and discrete diffusion baselines, our \method~revisits and selectively updates action tokens through token-level velocity modeling, thereby improving action consistency in robotic manipulation.

% Our approach explores two velocity-field constructions: an embedding-guided formulation that defines structured probability paths and an auxiliary velocity head formulation that predicts velocities directly from hidden states. It further integrates kinetic-optimal velocity design with a practical two-stage decoder that couples iterative refinement and deterministic validation, achieving superior execution quality and favorable inference efficiency across CALVIN, LIBERO, and real-world evaluations.

% \textbf{Future work.} 
% Future work includes scaling DFM-based VLA pretraining, improving  token refinement, and extending the framework to tighter joint modeling of perception, reasoning, and action in more open-ended real-world environments.





% ---- Bibliography ----
%
% BibTeX users should specify bibliography style 'splncs04'.
% References will then be sorted and formatted in the correct style.
%
